% =====================================================================
\chapter{Estimación puntual e intervalos de confianza}
\label{tema:4}
% =====================================================================

El \cref{tema:2} enseñó a juzgar un estimador; este tema enseña a construirlo. Se presentan los dos métodos generales de estimación puntual, el método de los momentos y el de máxima verosimilitud, y se aplican a los modelos de Bernoulli, Poisson, normal y exponencial. Después se pasa de la estimación puntual a la estimación por intervalos: intervalos de confianza para la media con varianza conocida y desconocida, para proporciones y cálculo del tamaño muestral. Se sigue el capítulo~9 de Wasserman~\cite{wasserman}, «Parametric Inference» (método de los momentos y máxima verosimilitud), las clases~12, 13 y~18 de Stats~200~\cite{stats200}, las lecturas~10b, 22 y~23a del curso 18.05 del MIT~\cite{mit1805} y el capítulo~8 de OpenStax~\cite{openstax}.

\section{El método de los momentos}

Sea $X_1,\dots,X_n$ una muestra \iid\ de un modelo $F_\theta$ con $\theta=(\theta_1,\dots,\theta_k)$. El \emph{momento poblacional} de orden $j$ es $m_j(\theta)=\E_\theta[X^{j}]$ y el \emph{momento muestral} de orden $j$ es $\widehat m_j=\frac1n\sum_iX_i^{j}$. Por la ley de los grandes números, $\widehat m_j\convp m_j(\theta)$.

\begin{definicion}[Estimador por el método de los momentos]\label{def:t4-momentos}
El \emph{estimador por el método de los momentos} $\widehat\theta$ es la solución del sistema de $k$ ecuaciones
\[
m_j(\widehat\theta)=\widehat m_j,\qquad j=1,\dots,k,
\]
es decir, el valor del parámetro con el que los primeros $k$ momentos del modelo coinciden con los de la muestra~\cite[cap.~9]{wasserman}.
\end{definicion}

\addtocounter{definicion}{1}%

El método es sencillo y produce estimadores consistentes en condiciones generales, pero no siempre los mejores y, a veces, incompatibles con los datos (\cref{eje:t4-uniforme-ab}).

\section{Máxima verosimilitud}

\begin{definicion}[Verosimilitud y estimador de máxima verosimilitud]\label{def:t4-emv}
Sean $x_1,\dots,x_n$ los datos observados y $f(x\mid\theta)$ la función de probabilidad (caso discreto) o la densidad (caso continuo) del modelo. La \emph{función de verosimilitud} y la \emph{log-verosimilitud} son
\[
\Lik(\theta)=\prod_{i=1}^{n}f(x_i\mid\theta),\qquad \ell(\theta)=\ln\Lik(\theta)=\sum_{i=1}^{n}\ln f(x_i\mid\theta).
\]
El \emph{estimador de máxima verosimilitud} (EMV) es el valor $\widehat\theta=\argmax_{\theta\in\Theta}\Lik(\theta)$~\cite[cap.~9]{wasserman}.
\end{definicion}

La verosimilitud es la probabilidad (o densidad) de los datos observados vista como función del parámetro, con los datos fijos. Como el logaritmo es estrictamente creciente, maximizar $\Lik$ equivale a maximizar $\ell$, que es una suma y se deriva con más facilidad. El procedimiento habitual es resolver la \emph{ecuación de verosimilitud} $\ell'(\theta)=0$ y comprobar que la solución es un máximo global.

\begin{teorema}[EMV de Bernoulli]\label{teo:t4-emv-bernoulli}
Si $x_1,\dots,x_n\in\{0,1\}$ son observaciones \iid\ de $\Ber(p)$ y $k=\sum_ix_i$ es el número de unos, con $0<k<n$, el EMV es $\widehat p=k/n=\xbar$.
\end{teorema}

Si $k=0$ o $k=n$, el máximo está en $\widehat p=0$ o $\widehat p=1$, respectivamente, y la fórmula $\widehat p=k/n$ sigue siendo válida.

\begin{teorema}[EMV de la normal]\label{teo:t4-emv-normal}
Si $x_1,\dots,x_n$ son observaciones \iid\ de $\Nor(\mu,\sigma^{2})$, no todas iguales, el EMV de $(\mu,\sigma^{2})$ es
\[
\widehat\mu=\xbar,\qquad\widehat\sigma^{2}=\frac1n\sum_{i=1}^{n}(x_i-\xbar)^{2}.
\]
\end{teorema}

El EMV de $\sigma^{2}$ coincide con el estimador por momentos y es sesgado (\cref{cor:t2-sesgo}); por eso, para datos, se usa casi siempre $S^{2}$.

\addtocounter{definicion}{1}%

\begin{observacion}[Propiedades del EMV]\label{obs:t4-propiedades}
Se enuncian sin demostración, como ampliación~\cite[cap.~9]{wasserman}: (a) \emph{invariancia}: si $\widehat\theta$ es el EMV de $\theta$, el EMV de $g(\theta)$ es $g(\widehat\theta)$; por ejemplo, el EMV de $\sigma$ es $\sqrt{\widehat\sigma^{2}}$; (b) en condiciones de regularidad (que no se cumplen en el \cref{ej:t4-uniforme}), el EMV es consistente y aproximadamente normal para $n$ grande, con un error estándar que se puede estimar a partir de la curvatura de $\ell$ en el máximo.
\end{observacion}

\section{Intervalos de confianza}

Una estimación puntual no informa de su precisión. Un intervalo de confianza acompaña la estimación de un margen de error con una garantía probabilística.

\begin{definicion}[Intervalo de confianza]\label{def:t4-ic}
Sea $0<\alpha<1$. Un \emph{intervalo de confianza} de nivel $1-\alpha$ para $\theta$ es un intervalo $[L,U]$ cuyos extremos $L=L(X_1,\dots,X_n)$ y $U=U(X_1,\dots,X_n)$ son estadísticos y tal que
\[
\PP_\theta(L\le\theta\le U)\ge1-\alpha\qquad\text{para todo }\theta\in\Theta.
\]
La probabilidad $\PP_\theta(L\le\theta\le U)$ se llama \emph{cobertura}~\cite[cap.~6]{wasserman}.
\end{definicion}

\begin{observacion}[Interpretación]\label{obs:t4-interpretacion}
Lo aleatorio es el intervalo, no el parámetro. Si se repitiera muchas veces el experimento y se calculara cada vez el intervalo de nivel $95\,\%$, aproximadamente el $95\,\%$ de esos intervalos contendría el verdadero $\theta$. Una vez observados los datos, el intervalo concreto (por ejemplo, $[0{,}79;\allowbreak\,0{,}84]$) contiene o no contiene a $\theta$; no es correcto decir que $\theta$ está en él «con probabilidad $0{,}95$» en el sentido frecuentista~\cite[lectura~23a]{mit1805}.
\end{observacion}

El método general para construir intervalos usa un \emph{pivote}: una función $Q(X_1,\dots,X_n;\theta)$ cuya distribución no depende de $\theta$.

\begin{teorema}[Intervalo para la media con $\sigma$ conocida]\label{teo:t4-ic-z}
Sean $X_1,\dots,X_n\sim\Nor(\mu,\sigma^{2})$ \iid\ con $\sigma$ conocida. El intervalo
\[
\Bigl[\Xbar-z_{1-\alpha/2}\frac{\sigma}{\sqrt n},\ \Xbar+z_{1-\alpha/2}\frac{\sigma}{\sqrt n}\Bigr]
\]
tiene cobertura exactamente $1-\alpha$. Si los datos no son normales pero $n$ es grande, la cobertura es aproximadamente $1-\alpha$.
\end{teorema}

La cantidad $E=z_{1-\alpha/2}\,\sigma/\sqrt n$ es el \emph{margen de error} (la semiamplitud del intervalo). Los valores más usados son $z_{0{,}95}=1{,}645$ (nivel $90\,\%$), $z_{0{,}975}=1{,}960$ ($95\,\%$) y $z_{0{,}995}=2{,}576$ ($99\,\%$).

\begin{corolario}[Tamaño muestral]\label{cor:t4-tamano}
Para que el margen de error del \cref{teo:t4-ic-z} sea a lo sumo $E$ basta tomar
\[
n\ge\Bigl(\frac{z_{1-\alpha/2}\,\sigma}{E}\Bigr)^{2},
\]
redondeando al entero superior.
\end{corolario}

\section{Intervalo para la media con \texorpdfstring{$\sigma$}{sigma} desconocida}

En la práctica $\sigma$ casi nunca se conoce. Al sustituirla por $S$ hay que reemplazar la normal por la $t$ de Student.

\begin{teorema}[Intervalo $t$]\label{teo:t4-ic-t}
Sean $X_1,\dots,X_n\sim\Nor(\mu,\sigma^{2})$ \iid\ con $\sigma$ desconocida y $n\ge2$. El intervalo
\[
\Bigl[\Xbar-t_{n-1;1-\alpha/2}\frac{S}{\sqrt n},\ \Xbar+t_{n-1;1-\alpha/2}\frac{S}{\sqrt n}\Bigr]
\]
tiene cobertura exactamente $1-\alpha$.
\end{teorema}

Como $t_{n-1;1-\alpha/2}>z_{1-\alpha/2}$, el intervalo $t$ es más ancho: refleja la incertidumbre adicional por estimar $\sigma$. Para $n$ grande (digamos $n\ge30$ o más) ambos cuantiles casi coinciden y, por el teorema central del límite, el intervalo es aproximadamente válido aunque los datos no sean normales.

\addtocounter{definicion}{1}%

\section{Intervalos para proporciones}

Si $X_1,\dots,X_n\sim\Ber(p)$, el \cref{cor:t3-binomial} da $\widehat p\approx\Nor\bigl(p,p(1-p)/n\bigr)$. El error estándar depende de $p$, desconocido, y se estima por $\sqrt{\widehat p(1-\widehat p)/n}$.

\begin{definicion}[Intervalo de Wald para una proporción]\label{def:t4-wald}
El intervalo de confianza aproximado de nivel $1-\alpha$ para $p$ es
\[
\widehat p\pm z_{1-\alpha/2}\sqrt{\frac{\widehat p(1-\widehat p)}{n}}.
\]
Como $p(1-p)\le1/4$, el intervalo \emph{conservador} $\widehat p\pm z_{1-\alpha/2}/(2\sqrt n)$ es siempre al menos tan ancho~\cite[§8.3]{openstax}.
\end{definicion}

La aproximación es buena si $n\widehat p$ y $n(1-\widehat p)$ son ambos al menos del orden de $10$; con muestras pequeñas o con $\widehat p$ cercano a $0$ o a $1$ la cobertura real es inferior a la nominal. Existen correcciones sencillas, como el método de «más cuatro» de~\cite[§8.3]{openstax} (sumar dos éxitos y dos fracasos antes de calcular) o el intervalo de Wilson, que quedan fuera de este módulo.

\begin{corolario}[Tamaño muestral para una proporción]\label{cor:t4-tamano-p}
Para que el margen de error sea a lo sumo $E$ sea cual sea $p$, basta $n\ge z_{1-\alpha/2}^{2}/(4E^{2})$. Si se dispone de una estimación previa $p^{*}$, basta $n\ge z_{1-\alpha/2}^{2}\,p^{*}(1-p^{*})/E^{2}$.
\end{corolario}

\addtocounter{definicion}{1}%

\addtocounter{definicion}{1}%

\addtocounter{definicion}{1}%

